\documentclass[11pt,a4paper]{article}
\usepackage[margin=2.5cm]{geometry}
\usepackage{booktabs}
\usepackage{longtable}
\usepackage{xcolor}
\usepackage[colorlinks=true,linkcolor=black,urlcolor=blue]{hyperref}
\usepackage{titlesec}
\usepackage{enumitem}
\usepackage{tcolorbox}
\titleformat{\section}{\Large\bfseries}{\thesection.}{0.6em}{}
\titleformat{\subsection}{\large\bfseries}{\thesubsection}{0.6em}{}
\titlespacing{\section}{0pt}{1.4em}{0.6em}
\setlength{\parskip}{0.55em}
\setlength{\parindent}{0pt}
\newcommand{\good}[1]{\textcolor{teal!70!black}{\textbf{#1}}}
\newcommand{\bad}[1]{\textcolor{red!65!black}{\textbf{#1}}}
\newtcolorbox{keybox}{colback=gray!7,colframe=gray!45,boxrule=0.5pt,left=8pt,right=8pt,top=6pt,bottom=6pt}

\title{\textbf{Mikhail LeTal --- where we are and how we win}\\[0.4em]
\large A plain account of what we have built, what happened tonight,\\and what the last 36 hours are for}
\author{Alexander Urvancev \and Yan Beletskiy}
\date{Wednesday 9 September 2026, 21:00 CEST}

\begin{document}
\maketitle
\thispagestyle{empty}

\begin{keybox}
\textbf{The one-paragraph version.} Our engine is uploaded, safe, and the strongest version we
have ever shipped. We lose games for one reason: our engine misjudges positions, so it walks
into losses believing it is winning. Tonight we tested our biggest idea for fixing that --- a
neural network evaluation --- and it lost decisively. We now have one machine, about 36 hours,
and room for exactly one more properly-measured experiment. Winning the whole event is not
realistic. A strong finish in Friday's tournament is, and this document says how.
\end{keybox}

\tableofcontents
\newpage

\section{Where we stand right now}

\subsection{What is live}

\textbf{Version 1.2 is uploaded and playing rated games.} It passed the safety gate this
evening: 500 games against test opponents, 500 wins, no crashes, no illegal moves, no losses on
time. That gate does not measure strength --- the opponents are trivial --- it measures whether
the engine can lose a game to something other than chess. It cannot.

\subsection{The three things happening tonight}

\begin{enumerate}[leftmargin=1.4em]
\item \textbf{A 1000-game test is running} on our move-ordering improvement. It finishes around
      06:20 Thursday. This is the last measurement that can change what we ship.
\item \textbf{The neural network is dead.} Tested tonight over 300 games, it lost 67 Elo. Details
      in Section 3.
\item \textbf{More training data is downloading} from the competition ladder --- 9{,}200 games,
      approved by you, running at one request per second.
\end{enumerate}

\subsection{The deadline, and what actually decides the winner}

\begin{center}
\begin{tabular}{ll}
\toprule
Friday 11 September, 11:00 London & Uploads close. Whatever is uploaded is locked. \\
Friday afternoon & 13-round Swiss tournament over the locked builds. \\
\bottomrule
\end{tabular}
\end{center}

The rules are explicit about what counts:

\begin{quote}
``\textbf{The ladder only seeds the final Swiss}''\\
``\textbf{Only the locked-build final Swiss counts, by points.}''\\
``Tie-breaks: points, Buchholz, head-to-head, \textbf{earlier final submission}.''
\end{quote}

\textbf{Our actual standing, from the platform itself:}

\begin{center}
\begin{tabular}{lllll}
\toprule
Rank & Rating & Record & Best streak & Checkmates \\
\midrule
\good{\#126 of 432} (top 29\%) & 1800 & 9--6--6 & 2 & 9 \\
\bottomrule
\end{tabular}
\end{center}

\textbf{That number decides almost nothing anyway.} It affects the first-round pairing on Friday
and nothing else. We joined late (our first rated game was round 67), started at a low rating,
and have therefore only ever been paired against weak opponents --- our opponents average well
below the field median and \textbf{we have never played a top-50 team}. \textbf{Thirteen games on
Friday decide everything.}

And note the last tie-break: \textbf{submitting earlier breaks ties}. That is a free half-point
in a close finish, and it costs us nothing.

\section{What we have built, and what each piece is worth}

Everything below was measured by playing hundreds of games against the previous version. Nothing
here is an estimate.

\subsection{The engine}

A chess engine written in Python, made fast with a compiler called numba. It searches roughly
800{,}000 positions per second and looks 10--14 moves ahead in a real game. \textbf{That part is
not our weakness.} For comparison, the version we started with searched 22{,}000 positions per
second.

\subsection{Version 1.2, the build now playing}

\begin{center}
\begin{tabular}{lr}
\toprule
What changed & Measured value \\
\midrule
A smarter search that avoids re-examining moves it already knows are bad & $+35$ Elo \\
Skipping positions already too good to need searching & $+15$ Elo \\
Three genuine bug fixes & --- \\
\midrule
\textbf{Total} & \good{$+50$ Elo} \\
\bottomrule
\end{tabular}
\end{center}

\textbf{The honest caveat:} the range of uncertainty on that $+50$ runs from about $-2$ to
$+102$. So the best estimate is $+50$, and the data cannot completely rule out that it is worth
nothing. It ships anyway, for reasons that are not the Elo number: one of the changes is
mathematically free (identical search, less work per position, so it cannot be worse), and one
of the bug fixes stops the engine scoring an already-won position as a draw --- which is a lost
game, not a lost rating point.

\subsection{How good is that, really?}

Elo comparisons across different opponents are unreliable. So we measured something better:
\textbf{average centipawn loss}, which asks a much simpler question.

\begin{keybox}
\textbf{What ``centipawn loss'' means.} A centipawn is 1/100th of a pawn --- the standard unit
engines use to score a position. For each move you played, we ask Stockfish (a far stronger
engine) how good the position was before your move and after it. \textbf{The drop is what your
move threw away.} Average that over every move.

A real example, from our round-85 loss: before our 13th move the position was scored $-24$
(roughly level). We played a move. After it, $-460$. That single move cost \textbf{436
centipawns} --- four and a half pawns --- in one decision.

This measure ignores ratings and pairings entirely. It compares move quality directly.
\end{keybox}

\begin{center}
\begin{tabular}{lrr}
\toprule
Team & Avg.\ centipawn loss & Moves losing $>$2 pawns \\
\midrule
Lubina            & \good{4.8}  & \good{0.0\%} \\
Emile Andrieu     & 9.3         & 0.4\% \\
Sobriety          & 16.7        & 0.2\% \\
THE ROOOOOKKK!!!! & 17.6        & 0.9\% \\
AI Fellows        & 18.8        & 0.7\% \\
\midrule
\textbf{Us, all builds}      & \bad{26.7} & \bad{2.4\%} \\
\textbf{Us, v1.1/v1.2 only}  & \bad{28.0} & \bad{1.8\%} \\
\bottomrule
\end{tabular}
\end{center}

\textbf{Two readings, and both matter.}

\good{The good news:} our blunder rate has nearly halved as we have improved the engine ---
3.1\% in the old builds, 1.8\% now. The work is landing.

\bad{The bad news:} we still throw away a serious amount of material \textbf{twice as often as
the worst team in the top five}. That single number is the entire gap between us and them. We do
not lose because our moves are slightly worse on average. We lose because roughly twice a game
we hand over a piece or a lost position.

\section{The core problem: the engine does not know when it is losing}

This is the most important section in the document. Everything else follows from it.

\subsection{What we found in two real games}

\textbf{Round 85 (lost by checkmate).} We traced every position against Stockfish. From move 14
to move 36 --- \textbf{23 consecutive moves} --- our engine believed it was winning by one to two
pawns. It was actually losing by two to six. The sign was wrong on 21 of those 23 moves.

It was not a blunder in the usual sense, and it was not time trouble. Move 13 was searched
\textbf{11 moves deep, over 812{,}000 positions}, and returned ``+45, I am fine.'' The position
it entered was scored $-460$. Eleven moves of lookahead did not reveal a five-pawn error.

\textbf{Round 87 (drawn, the first game v1.2 played).} Same disease, but this time the
evaluation was wrong \emph{in both directions in the same game}: too pessimistic by 100--370
centipawns for twenty moves, then too optimistic by 130--340 for the next ten.

\textbf{That second game is the more important one}, because a consistent bias could be corrected
with a constant. Noise cannot.

\subsection{Why more computing power cannot fix this}

\begin{keybox}
\textbf{The search is only as good as the judgement at the end of it.}

The engine looks ahead, reaches positions it cannot look past, and asks its evaluation:
``is this good for me?'' Every decision it makes traces back to that answer.

Our evaluation counts material and where the pieces sit. That is essentially all. It has
\textbf{no concept of activity at all} --- no measure of how many squares your pieces control.
So when it reaches a position where we are a pawn up, but the opponent has two bishops raking
open diagonals, a knight heading to a dominant square, and open lines at our king, it reports
\textbf{``+200, winning.''}

Searching deeper does not fix that. It just explores more thoroughly \emph{toward the wrong
goal.} You find the very best route to a position you have misjudged.
\end{keybox}

This is also why a faster development machine cannot help our playing strength: the competition
runs every engine on the same fixed hardware --- one CPU core, 2\,GB of memory. A better machine
helps us \emph{test} faster, which is worth something, but it adds no strength on Friday.

\section{What happened tonight: the network}

\subsection{The idea, and the result}

Instead of a hand-written evaluation, train a small neural network on a million positions and let
it learn to judge them. This was our biggest bet --- the only change with enough potential to
close the gap to the leaders.

\begin{center}
\begin{tabular}{lr}
\toprule
Network vs.\ v1.2, 300 games at the full 2-minute time control & \\
\midrule
Result & $+111$ wins, $21$ draws, $168$ losses \\
Strength & \bad{$-67$ Elo} (range $-106$ to $-29$) \\
\bottomrule
\end{tabular}
\end{center}

It is clearly, unambiguously worse. It does not ship.

\subsection{It is not the speed cost}

The obvious escape is ``the network is slower, so buy the speed back.'' The arithmetic says no.
The network costs 14.6\% of our search speed, which is about a tenth of a move of lookahead,
worth 5--9 Elo. The loss is 67. \textbf{So around 60 Elo of it is the network genuinely choosing
worse moves.}

\subsection{The strangest part: every measurement said it would win}

Before the games, we checked the network four separate ways. \textbf{It won all four.}

\begin{center}
\begin{tabular}{lll}
\toprule
Check & What it measured & Winner \\
\midrule
Held-out accuracy & Predicting Stockfish's score & network \\
Round-85 positions & Knowing who was winning, on a real loss & network (20/26 vs 5/26) \\
Quiet positions only & Accuracy where it should be strongest & network \\
\textbf{Game outcomes} & \textbf{Does the side it calls better actually win?} & \textbf{network} \\
\bottomrule
\end{tabular}
\end{center}

That last one was built specifically to be trustworthy --- no engine anywhere in the target, just
``did this side go on to win the game,'' measured over 12{,}000 positions from 9{,}200 real games.
The network scored 79.0\% against our hand-written evaluation's 76.3\%.

\textbf{And it lost 67 Elo.} Four independent measurements, all confidently wrong.

\section{Why the network failed --- you were right that it should be better}

A properly built neural evaluation is worth several hundred Elo in serious engines. Ours lost 67.
That is not evidence the approach fails. It is evidence we built it wrong, in six specific and
nameable ways.

\paragraph{1. Our network structurally cannot see king safety.}
Its input is 768 numbers: 12 piece types $\times$ 64 squares. \textbf{There is no king context
anywhere in it.} Serious networks use inputs of the form (\emph{our king's square}, piece,
square) --- around 41{,}000 inputs. That is what lets a network learn ``a knight on f5 is worth a
great deal \emph{because} the enemy king is on g8.''

Ours cannot learn that at all. And king safety is precisely the blindness that lost us round 85.
\textbf{We built a replacement that shares the original's blind spot.} No amount of training
fixes this one.

\paragraph{2. We asked it to predict something it cannot compute.}
We labelled every training position with Stockfish searching \emph{12 moves deep}. That score
contains tactics twelve moves away. Our network looks at a position and answers instantly, with
no search at all. For any position whose value depends on a tactic, \textbf{the correct answer is
unlearnable} --- and 98{,}700 parameters will happily memorise that unlearnable part as noise.
That is overfitting by construction.

\paragraph{3. Three to four orders of magnitude too little data.}
We trained on 186{,}000 positions. Serious training uses 100 million to a billion. Worse, ours
came from 3{,}356 games at about 51 positions each, so consecutive examples are highly similar
and the true independent sample is far smaller than it looks.

\paragraph{4. The wrong learning target.}
We minimised error on raw centipawn scores. Being 200 centipawns wrong in a level position is
catastrophic; being 200 wrong when you are already winning by 1400 is irrelevant. So most of the
learning was spent on positions that were already decided. Standard practice trains on
\emph{win probability} instead.

\paragraph{5. Trained on the wrong positions.}
We used positions from other engines' games. But the network is asked about positions
\emph{deep inside our own search} --- see Section 6, which turned out to be the most important
finding of the night.

\paragraph{6. Trained to answer the wrong question.}
The deepest one. We trained a function to \emph{predict what a search would conclude}. What the
engine needs is a function that \emph{correctly ranks the moves available right now}. Those are
different jobs, and nothing in our pipeline optimised the second. A network can be excellent at
the first and useless at the second --- which is exactly what happened, and exactly what four
misleading measurements were telling us.

\section{The most important thing we learned tonight}

\subsection{Both evaluations are terrible where it actually counts}

We captured 1{,}500 positions from \emph{inside a real search} --- the exact positions the engine
asks its evaluation about --- and compared them to positions from actual games.

\begin{center}
\begin{tabular}{lrr}
\toprule
Median error vs.\ Stockfish & Positions from games & \textbf{Positions inside the search} \\
\midrule
Our hand-written evaluation & 104 & \bad{342} \\
The network                 & 76  & \bad{304} \\
\bottomrule
\end{tabular}
\end{center}

\begin{keybox}
\textbf{Both evaluations are three to four times worse at exactly the positions where the engine
makes its decisions.}

And every measurement anyone ran today --- all four network checks, both game traces --- was done
on positions from games. Which are the positions a chess evaluation is \emph{least} often asked
about. We have been grading the engine on the wrong exam.
\end{keybox}

\subsection{And a second, related finding about the engine we are shipping}

The same data showed that at those search positions, Stockfish's typical score is around
\textbf{600 centipawns} where our evaluation reports \textbf{116}. Real positions deep in the
search are far more decisive than we think they are.

That matters because our search uses fixed thresholds --- ``if the position is at least 85
centipawns better than needed, stop searching this branch'' --- and \textbf{those thresholds were
chosen from textbooks, never fitted to our evaluation's actual scale.} It is a single constant,
with a plausible mechanism and no measurement in either direction. It is now a candidate for
Thursday.

\subsection{The lesson worth keeping}

Five separate defects today all had the same shape: a check that could not fail. A test set
defined using the training data. A ``wait until the machine is free'' loop whose own command
matched its own search pattern. A data filter that called the very function being tested. A guard
written to catch the incident that had already happened. A benchmark that verified how fast an
answer arrived but never that it was the same answer.

The question that catches all five, now written into our decision log:

\begin{center}
\emph{What is this check defined in terms of, and is that the thing I am testing?}
\end{center}

\section{The plan for the remaining hours}

\subsection{The hard constraint nobody can get around}

A properly-powered test --- 1000 games at the real time control --- takes \textbf{9.6 hours} on
our one machine. Thursday has roughly 06:20 to 22:00, minus time to freeze, safety-test and
upload.

\begin{keybox}
\textbf{Thursday holds exactly one properly-measured experiment.} Not two. Anything we squeeze in
at 300 games produces a result too vague to act on --- which is worse than no result, because a
vague number in a results table reads like evidence.
\end{keybox}

\subsection{The candidates for that one slot}

\begin{center}
\begin{longtable}{p{3.6cm}p{10.4cm}}
\toprule
\textbf{Mobility} & Adds a measure of piece activity to the evaluation --- the thing it is
completely missing, and the direct cause of the round-85 and round-87 misjudgements. Costs about
1.7$\times$ the time to reach a given depth, which is a lot. One unresolved test failure. \\
\addlinespace
\textbf{Threshold scaling} & The one-constant change from Section 6.2. Cheap to make, plausible
mechanism, entirely unmeasured. Not free: it also shifts reported scores, which touches how the
engine handles forced mates and manages its clock. \\
\addlinespace
\textbf{Neither} & Ship what we have, and spend Thursday on the freeze, the safety gate, an early
upload for the tie-break, and the report. \\
\bottomrule
\end{longtable}
\end{center}

\subsection{The decision: hold the slot}

\textbf{Threshold scaling is out, and the reason is worth understanding} --- it was my own
preference and the argument against it is decisive.

The two numbers point in opposite directions. ``Stockfish says 599 where we say 116'' describes
how much our evaluation \emph{moves}. ``Median error 304--342'' describes how \emph{wrong} it is.
Fitting a margin to the first says make it \textbf{smaller}; fitting it to the second says make it
\textbf{larger}. \textbf{The evidence does not tell us which way to turn the constant.} A
9.6-hour test on a coin flip dressed as a mechanism is the worst possible use of the only slot we
have.

There is a striking observation buried in that, though, worth recording for after the event:
\textbf{our pruning margins are 150--750 centipawns, and our own measured error at those
positions is 304--342.} Several of our margins sit below our own noise floor --- we are making
decisions on differences smaller than our uncertainty. And every pruning technique that consults
the static evaluation has measured badly on this engine. That is a coherent story, and it can be
tested cheaply without a screen: at the nodes where pruning fires, how often was the discarded
move actually the best one? Minutes, not hours, and it can run alongside anything.

\textbf{So: hold the slot rather than allocate it now.} Thursday's certain work --- freeze, safety
gate, early upload, report --- is worth more than an uncertain screen. If tonight's result lands
well and the freeze goes smoothly, mobility takes the slot. If not, we ship what we have.

\textbf{And mobility is not the safe bet the round-85 story makes it feel like.} It costs about
1.7$\times$ the time to reach a given depth, which is a $-10$ to $-23$ Elo handicap the new
evaluation term must overcome \emph{before} it is worth anything. The blindness it targets is real
and large; whether one term clears that bar is genuinely unknown.

A second network attempt is \textbf{not} on this list. It would cost about 7 hours to build plus
9.6 to test, which does not fit --- and after four measurements pointed the wrong way, there is
no cheap way to check our work as we go. The ceiling is high; the odds of hitting it blind, in
one attempt, are not good.

\subsection{Timetable}

\begin{center}
\begin{tabular}{lll}
\toprule
Tonight & 1000-game ordering test & finishes $\sim$06:20 \\
Thursday morning & Decide the one slot; start it & $\sim$9.6\,h \\
Thursday evening & Freeze, 500-game safety gate, \textbf{upload} & $\sim$1\,h \\
Friday morning & Re-upload only if something is clearly broken & \\
Friday 11:00 & Locked & \\
Friday afternoon & 13-round Swiss & \\
\bottomrule
\end{tabular}
\end{center}

\section{How to maximise our chances}

\subsection{Free points --- do all of these}

These cost essentially no machine time and depend on no measurement.

\begin{enumerate}[leftmargin=1.4em]
\item \textbf{Upload the final build early on Friday, not at 11:00.} ``Earlier final submission''
      is an explicit tie-break, and ties on points are common in a 13-round Swiss.
\item \textbf{Never lose a game to a crash, an illegal move, or the clock.} Each one is a whole
      point given away for nothing. Our 500-game safety gate exists exactly for this and must be
      run on whatever we finally ship.
\item \textbf{Keep uploading between now and the lock.} Ten uploads a day are allowed, the latest
      valid one plays, and a failed validation costs nothing. Every rated game is a free test of
      our own reliability.
\item \textbf{Fix the time allocation.} In round 87 the engine spent \textbf{1.0 second on the
      losing move while holding 88 seconds}, and 9.4 seconds on a move in an already-lost
      endgame. It finished with 32.9 seconds unused. Spending time in inverse proportion to how
      much it matters is a defect, it is cheap to investigate, and fixing it costs no strength.
\end{enumerate}

\subsection{The honest assessment}

\textbf{Winning the whole event is not realistic.} The top of the field is measurably better at
chess than we are --- a blunder rate of 0.0--0.9\% against our 1.8\% --- and closing that is an
evaluation rebuild, not a day of tuning. We spent our shot at one tonight.

\textbf{A strong Swiss finish is realistic, and is the right target.} Thirteen rounds decided on
points, over locked builds, with our ladder seeding nearly irrelevant. Consistency matters more
there than peak strength --- and consistency is exactly what a lower blunder rate buys. Ours has
already halved once.

\textbf{What would have to be true for a top-five finish:}

\begin{enumerate}[leftmargin=1.4em]
\item the move-ordering test lands positive tonight (likely --- expected $+15$ to $+25$);
\item Thursday's one experiment cuts the blunder rate meaningfully (uncertain);
\item we give away nothing to crashes, flags or lost time (entirely in our control);
\item the Friday pairings are not brutal in the early rounds (not in our control).
\end{enumerate}

\subsection{And the thing that is genuinely ours}

Every number in this document has a matched comparison, a stated range of uncertainty, and a
caveat attached wherever one is owed. Four times today a persuasive measurement said we were
winning and the games said otherwise --- and each time it was caught, because someone had written
down in advance what would count as evidence.

That is what we can defend in front of a judge, and it is worth more than any single technique.

\end{document}
